\section*{Discussion}

The assumption that improved drug representations would raise drug-blind prediction ceilings has driven a decade of encoder development in cancer pharmacogenomics. Our results redefine the constraint. The bottleneck is distributional, not representational. Mechanistically heterogeneous co-training suppresses the pathway-specific signals that any encoder would transmit regardless of its expressiveness, and restructuring the training distribution removes this suppression directly. Independent work has converged on the same null in drug-blind evaluation\cite{herbert2026monotherapy,bernett2025dreval} and in the drug-seen regime\cite{carli2025cellhit}, indicating that encoder development has been addressing the wrong variable.

The mechanism traces to training distribution heterogeneity. Co-training on mechanistically diverse drugs forces the model to learn a response map averaged across MoAs. Pathway-specific cell features that predict sensitivity for EGFR inhibitors conflict with those for topoisomerase inhibitors, and the model resolves this conflict by learning features predictive only of mean potency. A drug encoder, regardless of its expressiveness, feeds into a model whose response map has been smoothed across all MoAs. LINCS L1000 signatures capture a drug's average transcriptional effect on a reference population, not cell-specific response variation. Drug-target vectors are too coarse to resolve the pathway context differences that distinguish sensitive from resistant cell lines. The sole partial exception, MoA categorical identity in the Transformer, was abolished by label permutation, confirming that the residual signal reflects implicit class-conditional partitioning of the training distribution rather than representational content.

Within-MoA training distinguishes artifactual ceilings from genuine biological limits. Aligned training exceeded within-class profile concordances for ERK MAPK and EGFR, indicating that supervised models extract a generalizable mapping beyond profile similarity. Classes with high within-class mechanistic heterogeneity, such as Apoptosis regulation, are not addressable by distribution alignment alone because the training drugs within the class do not share a coherent response pattern to transfer. Response matching provides a complementary route, generalizing to patient-derived data (BeatAML) and to cell lines absent from training (CTRPv2), but requiring functional drug analogs in the reference set.

Checkpoint selection artifacts explain a substantial portion of reported progress as methodological inflation, consistent with the 72\% leakage rate documented across the field\cite{asiaee2026widespread}. Per-drug $r$, the cell-mean prior as a free baseline, and replicate concordance as a ceiling give the field evaluation tools to separate genuine advances from metric optimization.

Both approaches have scope conditions. Response matching requires functional drug analogs in the reference set, and within-MoA training requires mechanistically coherent drug classes with strong pathway dependence. Beyond these scope conditions, all results rest on cell-line IC$_{50}$ values, which abstract away \textit{in vivo} pharmacokinetics, tumor microenvironment effects, clonal heterogeneity, and drug combinations. Whether distribution alignment translates to improved patient response requires prospective validation in clinical cohorts.

The representation--distribution dissociation may extend beyond cancer pharmacogenomics. Any co-training distribution with mechanistic heterogeneity can suppress signals relevant to a specific test condition. Benchmarking should evaluate training distribution alongside representation quality to distinguish improved encoding from improved data curation.
